% lem-local: limit-shapes.tex:396-412 (label lem:local)
\paragraph{Paper excerpt.}
\begin{verbatim}
\begin{lemma}\label{lem:local}
Let $p>0$. There exists $C(d,\drift,\gauge,p)<\infty$ such that, for every integer~$n\geq2$, the following estimates hold simultaneously with probability at least $1-Cn^{-p}$.
\begin{enumerate}[label=\textup{(\roman*)}]
\item \underline{\emph{Largest local time}}:
\begin{equation}\label{eq:M}
\max_{x\in\Z^d}\ell_n(x)\leq C|A_n|^{\nf{1}{d}}+C(\log n)^2\, .
\end{equation}
\item \underline{\emph{Local times on time intervals}}: for all integers~$0\leq s<t\leq n$,
\begin{equation}\label{eq:interval}
M_{s,t}\leq Ck_{s,t}^{\nf{1}{d}}+C(\log n)^2\, .
\end{equation}
\item \underline{\emph{Approximation by the potential}}: for every $y\in\R^d$ with $|y|\leq2n$,
\begin{equation}\label{eq:approx}
|\widetilde\ell_n(y)-U_{D_n}(y)|\leq C\log n+C\begin{cases}\bigl(\max_x\ell_n(x)\bigr)^{\nf12}\log n,&d=2,\\\bigl(\max_x\ell_n(x)\log n\bigr)^{\nf12},&d\geq3\end{cases}\, .
\end{equation}
\end{enumerate}
\end{lemma}
\end{verbatim}
\paragraph{Lean statement.}
\begin{verbatim}
theorem CERW.Frozen.norm_local_time_potential {d : ℕ} (hd : 2 ≤ d) :
    ∀ Ψ : EuclideanSpace ℝ (Fin d) → ℝ, CERW.IsNorm Ψ →
    ∀ ε : ℝ, 0 < ε → (∀ i : Fin d, ε * Ψ (CERW.coordVec i) < 1 / (d : ℝ)) →
    ∀ p : ℝ, 0 < p → ∃ C : ℝ, 0 < C ∧
      ∀ ξ : Site d → EuclideanSpace ℝ (Fin d),
        (∀ x : Site d, x ≠ 0 → CERW.IsSubgradient Ψ (CERW.toSpace x) (ξ x)) → ξ 0 = 0 →
      ∀ {Ω : Type u} [MeasurableSpace Ω] (μ : Measure Ω) [IsProbabilityMeasure μ]
        (X : ℕ → Ω → Site d), CERW.IsDriftCERW μ ε ξ X → ∀ n : ℕ, 2 ≤ n →
        μ {ω | ¬ ((CERW.maxLocalTime (X · ω) n : ℝ)
                    ≤ C * ((CERW.departureRange (X · ω) n).card : ℝ) ^ ((1 : ℝ) / d)
                      + C * Real.log n ^ 2 ∧
                  (∀ s t : ℕ, s < t → t ≤ n →
                    (CERW.intervalMax (X · ω) s t : ℝ)
                      ≤ C * (CERW.freshCount (X · ω) s t : ℝ) ^ ((1 : ℝ) / d)
                        + C * Real.log n ^ 2) ∧
                  ∀ y : EuclideanSpace ℝ (Fin d), ‖y‖ ≤ 2 * n →
                    |CERW.cellLocalTime (X · ω) n y
                        - CERW.normPotential d ε Ψ (CERW.cellSet (X · ω) n) y|
                      ≤ C * Real.log n + C *
                        (if d = 2 then Real.sqrt (CERW.maxLocalTime (X · ω) n) * Real.log n
                          else Real.sqrt (CERW.maxLocalTime (X · ω) n * Real.log n)))}
          ≤ ENNReal.ofReal (C * (n : ℝ) ^ (-p))
--
\end{verbatim}
\paragraph{Transcription.}
Paper macros: \verb!\drift! is $\kappa$ (Lean \verb!ε!), \verb!\gauge! is $\Psi$, \verb!\nf{a}{b}! is $a/b$. \emph{Object mapping.} \begin{itemize} \item $\Z^d$, $\R^d$, $\gauge$, $\xi$ and the ellipticity condition are as in thm-norm-shape: \verb!Ψ! with \verb!CERW.IsNorm Ψ!; \verb!ε! with \verb!0 < ε! and \verb!∀ i, ε * Ψ (CERW.coordVec i) < 1 / d!; \verb!ξ! with \verb!∀ x, x ≠ 0 → CERW.IsSubgradient Ψ (CERW.toSpace x) (ξ x)! and \verb!ξ 0 = 0!; \verb!hd : 2 ≤ d!. The dependence of the paper's $C(d,\kappa,\gauge,p)$ is exactly the order of these binders. \item The walk with norm $\gauge$: \verb!X : ℕ → Ω → Site d! on an arbitrary probability space (\verb!Ω : Type u!, \verb![MeasurableSpace Ω]!, \verb!μ : Measure Ω!, \verb![IsProbabilityMeasure μ]!) with \verb!hX : CERW.IsDriftCERW μ ε ξ X!. This says: each \verb!X n! is measurable, \verb!X 0 = 0! almost surely, and every cylinder
